\section*{Limitations}

\paragraph{Platform scope.} The system requires Apple-Silicon hardware (MLX)
and Python 3.14. We argue this is a strength for the large Mac-owning cohort
(\S3)---readable pure-MLX modeling, unified memory, an ecosystem to graduate
into---but it does exclude learners on other platforms; a CUDA/CPU port is
the most requested extension path.

\paragraph{Scale.} All models are intentionally tiny (millions, not billions,
of parameters). The mechanisms are the real ones, but phenomena that emerge
only at scale (training instabilities, emergent capabilities, large-batch
serving dynamics) are described, not experienced. Similarly, the vision
module ships the config, tiling, and connector with a bit-identical text
path, but not a pretrained encoder or the multimodal training stages.

\paragraph{Pedagogical evaluation.} Verification in this paper is mechanical
(tests, drift guards, reproducible measurements). A controlled study of
learning outcomes is the principal item of future work, alongside classroom
deployment and comparison against slice-specific resources; the build labs
grade isolated mechanisms rather than in-place modifications of the live
model.

\paragraph{Coverage depth.} Some systems mechanisms are taught via their core
rule (e.g., the paged-KV address translation, the batching cohort key) rather
than the full threaded implementation; the deepest labs (ragged attention
masks, the complete serving loop) remain open contributions.

\paragraph{Training data.} The bundled runs use a small local code corpus for
convenience; the repository ships no dataset, and reproducing the reference
run requires the user to supply their own clearly licensed corpus.
